% Generated draft: requires mathematical and source review.
\documentclass[11pt]{article}
\usepackage[margin=25mm]{geometry}
\usepackage{fontspec}
\setmainfont{Latin Modern Roman}
\usepackage{amsmath,amssymb,mathtools,unicode-math}
\setmathfont{Latin Modern Math}
\usepackage{enumitem,needspace,xurl,hyperref}
\setlength{\parindent}{0pt}
\setlength{\parskip}{5pt}
\setlength{\emergencystretch}{3em}
\newcommand{\N}{\mathbb{N}}
\newcommand{\Z}{\mathbb{Z}}
\newcommand{\Q}{\mathbb{Q}}
\newcommand{\R}{\mathbb{R}}
\newcommand{\C}{\mathbb{C}}
\newcommand{\F}{\mathbb{F}}
\newcommand{\A}{\mathbb{A}}
\newcommand{\PP}{\mathbb{P}}
\newcommand{\E}{\mathbb{E}}
\newcommand{\eps}{\varepsilon}
\newcommand{\dd}{\,\mathrm{d}}
\begin{document}
\section*{UnsolvedMath ID: 30004231}
\textbf{Source problem code:} catalogue entry
\subsection{Definitions and mathematical statement}
Let \(P_1,P_2\subseteq\R^n\) be polytopes. Their Minkowski sum is
\[
P_1+P_2=\{x_1+x_2:x_1\in P_1,\ x_2\in P_2\}.
\]
The supplied source calls a face \(F=F_1+F_2\) totally mixed when
\[
F_i\subseteq P_i,\qquad \dim F_i>0\quad(i=1,2),\qquad
\dim F=\dim F_1+\dim F_2.
\]
It does not explicitly specify whether the summands \(F_i\) must be faces exposed by a common linear functional, or whether the whole polytope is included among the faces under consideration.

Under the conventional common-supporting-functional interpretation, for a nonempty polytope \(P\) and \(u\in\R^n\), write
\[
P^u=\{x\in P:\langle u,x\rangle=\max_{y\in P}\langle u,y\rangle\}.
\]
Then \((P_1+P_2)^u=P_1^u+P_2^u\). A precise candidate for the totally mixed face poset, including the whole polytope whenever it satisfies the criterion, is
\[
\mathcal T(P_1,P_2)=\left\{(P_1+P_2)^u:
\begin{array}{l}
u\in\R^n,\ \dim P_1^u>0,\ \dim P_2^u>0,\\
\dim (P_1+P_2)^u=\dim P_1^u+\dim P_2^u
\end{array}\right\},
\]
ordered by inclusion. This is an explicit interpretive convention, requiring confirmation against the intended source formulation. If only proper faces are intended, the whole polytope must instead be excluded.

For a finite poset \(\mathcal T\), its order complex \(\Delta(\mathcal T)\) has the elements of \(\mathcal T\) as vertices and its nonempty finite chains as simplices. Interpret the geometric realization of the poset as \(|\Delta(\mathcal T)|\), and shellability of the poset as shellability of this order complex.

With these qualifications, the supplied question asks whether, for every pair of polytopes in the intended class, \(\Delta(\mathcal T(P_1,P_2))\) is shellable and its realization has the rational homology of a finite wedge of \((n-2)\)-spheres. For \(n\geq2\), the latter assertion means that there exists an integer \(r\geq0\), depending on the pair, such that
\[
H_*\bigl(|\Delta(\mathcal T(P_1,P_2))|;\Q\bigr)
\cong
H_*\left(\bigvee_{j=1}^{r}S^{n-2};\Q\right)
\]
as graded rational vector spaces, with the zero-fold wedge interpreted as a point. The source supplies no lower bound on \(n\), no full-dimensionality or general-position hypothesis, and no convention for an empty totally mixed face poset. It also does not specify whether shellability means pure shellability or a version allowing nonpure complexes. These issues remain unresolved; the displayed formulation is conditional on the stated conventions.
\subsection{Short English statement}
For two polytopes, consider faces of their Minkowski sum whose two summands have positive dimension and whose dimensions add. The question asks whether their inclusion poset is shellable and whether its order complex has the rational homology of a wedge of \((n-2)\)-spheres. The intended summands, face conventions, admissible dimensions, and shellability convention require clarification.
\subsection{Sources}
\begin{itemize}
\item Local UnsolvedMath registry, supplied complete source record: permanent ID \(30004231\), title \emph{Shellability and Homology of Totally Mixed Face Posets}, exact source problem code \texttt{catalogue entry}. No individual proposer is supplied; the proposed year is recorded as 2019.
\item Supplied report citation: \emph{Geometric, Algebraic, and Topological Combinatorics} (2020), Oberwolfach Reports, pp.~2395--2472, DOI \texttt{10.4171/owr/2019/39}. \url{https://doi.org/10.4171/owr/2019/39}. The background attributes the problem to workshop open-problem discussions documented in Oberwolfach Reports.
\item Supplied literature source, labeled \emph{Mixed subdivisions and faces}: \url{https://arxiv.org/abs/1607.00347}. No authors or further bibliographic details are supplied.
\item Supplied literature assessment dated exactly \texttt{2026-08-22}: the record labels the problem open and states that results on mixed-face counts do not establish shellability or the claimed rational wedge-of-spheres homology for this totally mixed-face poset. The listed sources are the report and arXiv link above. This is supplied dated curation, described in the background as literature triage derived from the supplied OWR curation; neither its status nor its cited claims have been independently verified here.
\end{itemize}
\end{document}
